\documentclass[letter]{article}

\usepackage{fullpage}
\usepackage{url}
\usepackage{hyperref}


\usepackage{listings}
\usepackage{color}
\usepackage{graphics}
\usepackage{graphicx}

\newcommand{\question}{{\noindent \bf Question: }}
\newcommand{\todo}[1]{{\bf #1}}


\title{CUDA: Polynomial Expansion}
\date{}

\begin{document}

\maketitle

\section{Preliminary}

To use CUDA, add \texttt{module load cuda/11.2} in your \texttt{.bashrc}.% or
%\texttt{.bash\_profile}.  %Also, if you have a \texttt{module load gcc}
%in there, replace it with \texttt{module load gcc/5.3.0}

%To submit a CUDA job manually, use \texttt{-lnodes=1:ppn=7:gpus=1} in
%you \texttt{qsub} command. But \texttt{make bench\_gpu} will take care
%of that.

%% Processing on the GPU can be asynchronous, which can lead to time
%% measurement of 0 seconds. Use
%% \texttt{cudaDeviceSynchronize()} to ensure previously submitted tasks are
%% complete. (memcpy and kernel launch are tasks.)

\section{Polynomial expansion}

The problem is simple. It is to compute a polynomial function $F$ of degree $d$:

$$F(x) = \sum_{i=0}^d a_i x^i$$

for an array of $n$ values of x.

(The code for polynomial expansion on the CPU is provided.)

\question Write a simple CUDA code that allocates and fill an array on
the CPU and transfer it to the GPU. (Take array size as a parameter)

\question Compute the polynomial expansion of each element of the
array on the GPU. (Take block size and degree of the polynomial as a
parameter.)

\question Bring the results back on the CPU and confirm the GPU code
is correct. You can test your code by running \texttt{make test} which
will queue both CPU and GPU jobs to the cluster which will ouput
\texttt{cputest-\textit{xyz}.out} which contains the desired answer, and
\texttt{gputest-\textit{xyz}.out} which contain your gpu code answer.

Hint: Remember to check explicitly for errors!

\section{Measurements}

%% \question Measure PCI-express latency. (That is the time for an array of size 1.)

%% \question Measure PCI-express Bandwidth. (The initial memory copy with different sizes of the array.)

%% \question Measure GPU memory bandwidth. (Exclude memory copies and use the simplest polynomial.)

%% \question Measure GPU flops rating. (Exclude memory copies and use a complex polynomial.)

\question Measure the runtimes of the codes with \texttt{make
  bench} for polynomial function ofdegree 1, 10, 100, and 1000.

\question Compare to the performance of the CPU implementation. You
can compute a speedup table with \texttt{make table}, this will output
the time and speedup in the form of a text table in \texttt{resulttable.txt}.
Which is faster in which configuration? Why do you think that is?

\end{document}
